% ==================== 7 结论 ====================
\setcounter{section}{6}
\section{结　论}

本文把课程两次建模作业整理成文：先用 Diamond 建立了 FCC、BCC、HCP、NaCl、金红石、
锐钛矿、萤石与金刚石八种结构，再对六类对称元素做了仿真演示。下面按
``结构—方法—对称性''三个方面归纳所得。

八个结构覆盖了金属、离子晶体与共价晶体的主要类型。金属单质中，FCC 与 HCP 的堆积
因子同为 0.7405，区别只在密排层的堆积顺序（ABCABC 与 ABAB）；BCC 的堆积因子降到
0.6802、配位数为 8，却是许多高温相与碱金属的稳定结构。离子晶体的配位数由阴阳离子
半径比决定：NaCl 中 Na$^+$ 取八面体六配位；萤石中 Ca$^{2+}$ 为立方八配位、F$^-$ 为
四面体四配位。TiO$_2$ 的两种晶型同为 TiO$_6$ 八面体配位，差别集中在连接方式：金红石
共棱成链、链间共角，锐钛矿则全部共边。金刚石结构以 $sp^3$ 四面体配位为特征，堆积
因子仅 0.3401，与共价键的方向性密切相关。

方法上，建模的关键是让``空间群—晶格常数—非对称单元原子''三者自洽：参数设对之后，
晶胞填充、晶面显示与配位多面体都只是顺次操作；粉末衍射模拟则把结构模型和真实的
XRD 实验联系了起来。

对称性方面，六类元素的演示把抽象符号变成了可对照的图像：反演中心让原子成对、关于
中心等距反向（$P\,\bar1$ 单胞 2 个原子）；3 次旋转轴给出 3 个共面原子（$P\,3$）；
旋转反演轴 $\bar3$ 相当于 $3+\bar1$ 的联合操作，产生上下两层共 6 个原子，沿轴观察时
呈``似六重''的假象；$P\,m\,\bar3$ 中原子关于三组互相垂直的镜面镜像成对；$3_1$ 螺旋轴
与 $c$ 滑移面含平移分量，螺旋上升与错位排列只有在超晶胞中才显示出来。这些演示说明：
对称性既决定单胞内的原子分布，也决定衍射的系统消光规律。

经过这次作业，笔者对``空间群—原子占位—配位多面体—堆积方式''的分析路径有了具体的
体会，也熟悉了用对称符号反推结构信息的做法；这些模型与截图将直接用于后续 XRD 谱的
指标化对照。
